% !TITLE 临江仙·柳絮
% !AUTHOR 曹雪芹
% !DYNASTY 明清
% !GENRE 词
% !ORDER 7
% !SOURCE 《红楼梦》（借薛宝钗之手）
% !EDITION 《红楼梦》，中国艺术研究院红楼梦研究所校注，人民文学出版社（2008年版为候选）
% !EDITION_STATUS pending

\begin{poem}
白玉堂\textsuperscript{1}前春解舞\textsuperscript{2}，东风卷得均匀。
蜂围蝶阵\textsuperscript{3}乱纷纷。几曾随逝水\textsuperscript{4}？岂必委芳尘\textsuperscript{5}？

万缕千丝终不改，任他随聚随分。
韶华\textsuperscript{6}休笑本无根。好风凭借力，送我上青云\textsuperscript{7}。
\end{poem}

\begin{annotations}
\notetext{1}{白玉堂：白玉砌成的厅堂，指富贵人家的宅第。}
\notetext{2}{春解舞：春（风）懂得如何舞蹈。把柳絮在风中飞舞比作春在跳舞。}
\notetext{3}{蜂围蝶阵：蜜蜂围成圈、蝴蝶排成阵。形容春天蜂蝶纷飞的景象——也暗喻众人对柳絮的围观与嘲笑。}
\notetext{4}{逝水：一去不返的流水。谁说柳絮一定随水而去、不复存在？}
\notetext{5}{委芳尘：委身于尘土之中。委，弃置。谁说柳絮一定落入泥土化为尘埃？}
\notetext{6}{韶华：美好的春光，也指年轻的时光、青春。春光不要笑柳絮本来没有根——无根之物也可以凭借风力高飞。}
\notetext{7}{青云：高空的云，比喻高位、远大的前程。“青云直上”指仕途顺利。薛宝钗借柳絮自喻：不靠出身，凭自身力量，也要飞上青云。}
\end{annotations}

\begin{appreciation}
这阕《临江仙》出自《红楼梦》第七十回。暮春时节，史湘云偶填了一首咏柳絮的《如梦令》，众人来了兴致，各以柳絮为题填词。林黛玉那首《唐多令》写得凄清，嫁与东风春不管，凭尔去，忍淹留，一句话把柳絮说成了任人摆布的弃物。薛宝钗偏要反着写，众人看罢，都称这一首翻得好，推它为尊。

柳絮在古诗里一向不是好身份。它无根无蒂，风一吹便散，落在水里随水去落在泥里随泥化。谢道韫拿柳絮因风起比拟飞雪，取的是它的轻盈；苏轼写杨花，说细看来点点是离人泪，取的是它的飘零；杜甫更直截，颠狂柳絮随风舞，干脆说它轻狂。宝钗偏偏处处都不肯认这个命。

白玉堂前春解舞，东风卷得均匀——同样是一阵东风，别人写来是飘零，她偏说是起舞，而且舞得从容卷得均匀不急不乱。蜂围蝶阵乱纷纷是热闹，也是围观：众人都看着这无根之物，看它几时落下来。几曾随逝水？岂必委芳尘？两个反问顶回去——谁说柳絮注定要随流水漂走埋进尘土？这两句才是全词的关口，她不肯接受别人替它安排好的结局。

下阕一句顶一句。万缕千丝终不改，柳丝千条万条，细细的，却都是自己的，任凭聚散，心志不移；韶华休笑本无根，春光你笑我无根，无根又如何。最末一句点题：好风凭借力，送我上青云。不必靠门第出身只等一阵好风一样上得青云。

咏物从来是咏人。同是一片柳絮，黛玉看见身不由己的漂泊，宝钗看见可以借力的东风，这正是两个人各自的处境与心气。曹雪芹写这一首，也未必全是夸赞——一个把“青云”二字挂在嘴边的少女，日后究竟上没上去，读下去自然知道。不过敢写下“送我上青云”这五个字，已经见出她不肯认输的性情了。
\end{appreciation}
